\chapter{Produit scalaire et orthogonalité}
\section*{Exercice 16.1 -- Coordonnées et angle}
Pour $u=(1,2)$ et $v=(3,-1)$, calculer produit scalaire, normes et cosinus de l'angle.

\section*{Exercice 16.2 -- Orthogonalité}
Déterminer $t$ pour que $(t,2)$ soit orthogonal à $(3,-6)$.

\section*{Exercice 16.3 -- Identité de norme}
Démontrer $\|u+v\|^2+\|u-v\|^2=2\|u\|^2+2\|v\|^2$.

\section*{Exercice 16.4 -- Projection signée}
Si $\|u\|=4$, $\|v\|=3$ et l'angle entre eux vaut $2\pi/3$, calculer $u\cdot v$ et la projection algébrique de $v$ sur la direction de $u$.

\section*{Exercice 16.5 -- Cauchy--Schwarz}
Démontrer $(x+y+z)^2\le3(x^2+y^2+z^2)$ et préciser l'égalité.

\section*{Exercice 16.6 -- Base non orthonormée}
Deux vecteurs de base unitaires $e_1,e_2$ font un angle $\pi/3$. Calculer $\|e_1+e_2\|^2$ et expliquer pourquoi la somme des carrés des coordonnées ne convient pas.
